% Options for packages loaded elsewhere
% Options for packages loaded elsewhere
\PassOptionsToPackage{unicode}{hyperref}
\PassOptionsToPackage{hyphens}{url}
\PassOptionsToPackage{dvipsnames,svgnames,x11names}{xcolor}
%
\documentclass[
  letterpaper,
  DIV=11,
  numbers=noendperiod]{scrartcl}
\usepackage{xcolor}
\usepackage{amsmath,amssymb}
\setcounter{secnumdepth}{5}
\usepackage{iftex}
\ifPDFTeX
  \usepackage[T1]{fontenc}
  \usepackage[utf8]{inputenc}
  \usepackage{textcomp} % provide euro and other symbols
\else % if luatex or xetex
  \usepackage{unicode-math} % this also loads fontspec
  \defaultfontfeatures{Scale=MatchLowercase}
  \defaultfontfeatures[\rmfamily]{Ligatures=TeX,Scale=1}
\fi
\usepackage{lmodern}
\ifPDFTeX\else
  % xetex/luatex font selection
\fi
% Use upquote if available, for straight quotes in verbatim environments
\IfFileExists{upquote.sty}{\usepackage{upquote}}{}
\IfFileExists{microtype.sty}{% use microtype if available
  \usepackage[]{microtype}
  \UseMicrotypeSet[protrusion]{basicmath} % disable protrusion for tt fonts
}{}
\makeatletter
\@ifundefined{KOMAClassName}{% if non-KOMA class
  \IfFileExists{parskip.sty}{%
    \usepackage{parskip}
  }{% else
    \setlength{\parindent}{0pt}
    \setlength{\parskip}{6pt plus 2pt minus 1pt}}
}{% if KOMA class
  \KOMAoptions{parskip=half}}
\makeatother
% Make \paragraph and \subparagraph free-standing
\makeatletter
\ifx\paragraph\undefined\else
  \let\oldparagraph\paragraph
  \renewcommand{\paragraph}{
    \@ifstar
      \xxxParagraphStar
      \xxxParagraphNoStar
  }
  \newcommand{\xxxParagraphStar}[1]{\oldparagraph*{#1}\mbox{}}
  \newcommand{\xxxParagraphNoStar}[1]{\oldparagraph{#1}\mbox{}}
\fi
\ifx\subparagraph\undefined\else
  \let\oldsubparagraph\subparagraph
  \renewcommand{\subparagraph}{
    \@ifstar
      \xxxSubParagraphStar
      \xxxSubParagraphNoStar
  }
  \newcommand{\xxxSubParagraphStar}[1]{\oldsubparagraph*{#1}\mbox{}}
  \newcommand{\xxxSubParagraphNoStar}[1]{\oldsubparagraph{#1}\mbox{}}
\fi
\makeatother


\usepackage{longtable,booktabs,array}
\usepackage{calc} % for calculating minipage widths
% Correct order of tables after \paragraph or \subparagraph
\usepackage{etoolbox}
\makeatletter
\patchcmd\longtable{\par}{\if@noskipsec\mbox{}\fi\par}{}{}
\makeatother
% Allow footnotes in longtable head/foot
\IfFileExists{footnotehyper.sty}{\usepackage{footnotehyper}}{\usepackage{footnote}}
\makesavenoteenv{longtable}
\usepackage{graphicx}
\makeatletter
\newsavebox\pandoc@box
\newcommand*\pandocbounded[1]{% scales image to fit in text height/width
  \sbox\pandoc@box{#1}%
  \Gscale@div\@tempa{\textheight}{\dimexpr\ht\pandoc@box+\dp\pandoc@box\relax}%
  \Gscale@div\@tempb{\linewidth}{\wd\pandoc@box}%
  \ifdim\@tempb\p@<\@tempa\p@\let\@tempa\@tempb\fi% select the smaller of both
  \ifdim\@tempa\p@<\p@\scalebox{\@tempa}{\usebox\pandoc@box}%
  \else\usebox{\pandoc@box}%
  \fi%
}
% Set default figure placement to htbp
\def\fps@figure{htbp}
\makeatother





\setlength{\emergencystretch}{3em} % prevent overfull lines

\providecommand{\tightlist}{%
  \setlength{\itemsep}{0pt}\setlength{\parskip}{0pt}}



 


\KOMAoption{captions}{tableheading}
\makeatletter
\@ifpackageloaded{caption}{}{\usepackage{caption}}
\AtBeginDocument{%
\ifdefined\contentsname
  \renewcommand*\contentsname{Table of contents}
\else
  \newcommand\contentsname{Table of contents}
\fi
\ifdefined\listfigurename
  \renewcommand*\listfigurename{List of Figures}
\else
  \newcommand\listfigurename{List of Figures}
\fi
\ifdefined\listtablename
  \renewcommand*\listtablename{List of Tables}
\else
  \newcommand\listtablename{List of Tables}
\fi
\ifdefined\figurename
  \renewcommand*\figurename{Figure}
\else
  \newcommand\figurename{Figure}
\fi
\ifdefined\tablename
  \renewcommand*\tablename{Table}
\else
  \newcommand\tablename{Table}
\fi
}
\@ifpackageloaded{float}{}{\usepackage{float}}
\floatstyle{ruled}
\@ifundefined{c@chapter}{\newfloat{codelisting}{h}{lop}}{\newfloat{codelisting}{h}{lop}[chapter]}
\floatname{codelisting}{Listing}
\newcommand*\listoflistings{\listof{codelisting}{List of Listings}}
\makeatother
\makeatletter
\makeatother
\makeatletter
\@ifpackageloaded{caption}{}{\usepackage{caption}}
\@ifpackageloaded{subcaption}{}{\usepackage{subcaption}}
\makeatother
\usepackage{bookmark}
\IfFileExists{xurl.sty}{\usepackage{xurl}}{} % add URL line breaks if available
\urlstyle{same}
\hypersetup{
  pdftitle={Supplementary Methods},
  pdfauthor={Kwaku Peprah Adjei},
  colorlinks=true,
  linkcolor={blue},
  filecolor={Maroon},
  citecolor={Blue},
  urlcolor={Blue},
  pdfcreator={LaTeX via pandoc}}


\title{Supplementary Methods}
\author{Kwaku Peprah Adjei}
\date{Invalid Date}
\begin{document}
\maketitle

\renewcommand*\contentsname{Table of contents}
{
\hypersetup{linkcolor=}
\setcounter{tocdepth}{3}
\tableofcontents
}

\section{Introduction}\label{introduction}

This document provides the detailed methods for data processing,
covariate construction, statistical modelling, prediction, and
validation used in the habitat-selection analyses reported in the main
manuscript. We model the probability that a dive location is a ``used''
site (observed dives) versus an ``available'' site (random background)
as a function of distance to colony, proximity to harvest activity,
bathymetry, and other covariates. Models are fitted using flexible
generalized additive models (GAMs) with individual-level random effects
and a spatial smoother to account for residual spatial structure.

The goals of this methods description are to:

\begin{itemize}
\tightlist
\item
  Describe the data sources and the cleaning/processing steps used to
  produce the modelling dataset;
\item
  Define the covariates and how they were derived (including harvest
  timing variables);
\item
  Present the GAM formulation used for inference and prediction, with
  justification of smoothing and random-effect choices;
\item
  Explain prediction procedures and the production of maps and summary
  plots; and
\item
  Describe validation and reproducibility practices.
\end{itemize}

\section{Study area and data sources}\label{study-area-and-data-sources}

\subsection{Study area}\label{study-area}

Analyses were centred on a single breeding colony and restricted to
locations within a \(33\) km radius of the colony centroid. All spatial
calculations were performed in a projected coordinate reference system
with metre units to ensure consistent buffering and distance
computations.

\subsection{Primary data sources}\label{primary-data-sources}

\begin{itemize}
\item
  \textbf{Colony location}: single point used to calculate
  distance-to-nest (\texttt{DistNest}).
\item
  \textbf{Kelp forest polygons} (\texttt{tareskog}): used to define
  available habitat for diving seabirds.
\item
  \textbf{Dive locations}: point observations for two sampling years
  (\(2023\) and \(2024\)). Each dive record includes at minimum a unique
  individual identifier (\texttt{ID}), a timestamp
  (\texttt{UTC\_datetime}), sex, and logger model.
\item
  \textbf{Harvest events}: spatial records (points or polygons) with
  timestamps; used to infer harvest overlap and timing relative to dive
  observations.
\item
  \textbf{Environmental covariates}: GEBCO bathymetry (raster, NetCDF)
  and substrate shapefiles (hard rock, sediment).
\end{itemize}

All raw spatial layers were reprojected to a common metric CRS before
analysis.

\subsection{Data preparation and
cleaning}\label{data-preparation-and-cleaning}

This section outlines the preprocessing steps applied before statistical
modelling.

\subsubsection{Restrict observations to available
habitat}\label{restrict-observations-to-available-habitat}

Observations were retained only if they occurred within the mapped kelp
polygons (tareskog) clipped to the colony-centred 33 km buffer.
Restricting to mapped kelp ensures that predictions are constrained to
biologically plausible habitat where diving is possible.

\subsubsection{Resolving duplicate
intersections}\label{resolving-duplicate-intersections}

Spatial intersection of point data with polygonal kelp layers
occasionally created duplicate rows for a single dive when the point
overlapped multiple polygon parts. To remove these duplicates while
preserving observational information, we grouped by the dive identifier
and a geometry identifier (derived from the geometry text), and
summarised attributes by retaining the first non-missing value for each
variable; a single geometry per dive was kept. This approach reduces
artefactual duplication without discarding unique events.

\subsubsection{Temporal and spatial thinning
(optional)}\label{temporal-and-spatial-thinning-optional}

To reduce temporal clustering and spatial autocorrelation that can
inflate Type I error rates, an optional thinning procedure was applied
at the level of individual × day. Within each day for a given
individual, DBSCAN clustering was performed on projected coordinates
(eps set in metres, minPts = 1). From each spatial cluster the earliest
timestamped observation was retained. This thinning is configurable and
used to evaluate sensitivity to pseudo-replication.

\#\#\#Sampling available (background) points Available (background)
points were sampled to represent habitat available to animals under
plausible accessibility rules. Two strategies were used depending on
analysis goals:

\begin{itemize}
\item
  Individual-specific accessible area: for each individual, we created
  an accessible polygon by buffering that individual's observed points
  by a radius (ID\_buffer\_m) and intersecting the resulting buffer with
  the kelp mask. Background points were sampled uniformly at random from
  this per-individual area; sampling intensity was n\_per\_point
  background points per observed point.
\item
  Group-level accessible area: individuals were grouped by sex and
  logger model. For each group, we excluded small buffers around
  observed locations (min\_dist\_m) that represent known-use zones and
  sampled uniformly from the remainder of the study region intersected
  with the kelp mask.
\end{itemize}

Harvest-aware stratification (optional): when harvest maps were
available, we created a harvest buffer (e.g., 200 m) and stratified
background sampling so that at least a target fraction
(min\_frac\_in\_bg, e.g., 0.25) of background points fell inside the
harvest buffer. This stratification ensured sufficient representation of
both harvested and unharvested strata for modelling and inference. To
account for unequal sampling areas, background points sampled from
different strata were assigned area-based sampling weights (area of
stratum divided by number of points sampled from that stratum) and
weights were normalised to have mean 1 prior to modelling.

Time-matching: timestamps for available points were randomly sampled
from the set of observed timestamps for the same individual (or group)
to match temporal distributions of used and available points and reduce
confounding due to temporal variation in environmental or harvest
conditions.

\subsubsection{Harvest overlap and timing
covariates}\label{harvest-overlap-and-timing-covariates}

For each used and available point, we computed covariates summarising
harvest overlap and relative timing:

\begin{itemize}
\item
  Buffer overlap indicators: For each buffer radius r (e.g., r = 50 m
  and 200 m), we created a binary indicator indicating whether the point
  fell inside a harvest buffer of radius r.
\item
  Earliest overlapping harvest time: Using an overlay (spatial join)
  between points and harvest buffers, we recorded the earliest harvest
  time overlapping each point (if any). When a point overlapped multiple
  harvest events we retained the minimum (earliest) timestamp per point.
\item
  Temporal relation to harvest: We computed DiffTime\_r =
  harvest\_time\_r − observation\_time, where negative values indicate
  that the harvest occurred before the observation (i.e., the site was
  harvested prior to the dive), and positive values indicate harvest
  after the dive. Points were classified per radius r into categories:
  ``after'' (harvest before the observation), ``before'' (harvest after
  the observation), ``same\_time'' (harvest time equals observation
  time), and ``noharvest'' (no overlapping harvest). From these
  categories we derived binary covariates used in the GAM: harvested50m
  = 1 if category ``after'' at r = 50 m, else 0 (and similarly for
  harvested200m).
\end{itemize}

Rationale: these covariates allow estimation of short-term (local)
responses to harvest (e.g., avoidance or attraction) and examination of
whether effects differ by sex or period.

\subsection{Environmental covariates and prediction
grid}\label{environmental-covariates-and-prediction-grid}

\subsubsection{Raster predictors (10 m
grid)}\label{raster-predictors-10-m-grid}

Environmental predictors were prepared on a 10 m resolution grid
covering the study extent:

Bathymetry: GEBCO netCDF rasters were reprojected to the study CRS,
resampled to 10 m, and standardised (mean 0, SD 1).

Substrate: hard rock and sediment polygon layers were rasterised to
binary rasters (1 = presence of substrate class, 0 = absence).

Kelp mask: tareskog polygons were rasterised to a binary mask used to
constrain predictions to available habitat.

Distance rasters: for each year we computed distance-to-harvest rasters
(meters) and a distance-to-colony raster (meters). Distances were
standardised prior to modelling.

\subsubsection{Prediction grid}\label{prediction-grid}

We generated a regular grid of prediction points (grid centres) within
the kelp mask; grid spacing (e.g., 100 m for mapping or 10 m for
fine-scale maps) was chosen to balance spatial resolution and
computational cost. Predictions were computed at grid points and
converted to raster format for visualization and for spatial
aggregation.

\section{Statistical model: GAM formulation and
interpretation}\label{statistical-model-gam-formulation-and-interpretation}

We modelled the probability that a location is used (a dive) using a
binomial generalized additive model (GAM). GAMs allow flexible
non-linear relationships via smooth functions while retaining parametric
interaction terms for testing specific hypotheses.

\subsection{Model notation and
equation}\label{model-notation-and-equation}

Let \(i\) index observations (used or available). For observation \(i\),
let \(y_i ∈ \{0, 1\}\) denote the response (1 = used, 0 = available).
The model is:

\subsubsection{Response model (logistic
scale)}\label{response-model-logistic-scale}

\[logit(p_i) = η_i\]

where \(p_i = Pr(y_i = 1)\), and \(η_i\) is the linear predictor.

\subsubsection{Linear predictor}\label{linear-predictor}

\[
\begin{split}
η_i &= f1(DistNest_i) + β1 × DistHarvest_i × Sex_i +  \\
& β2 × DistHarvest_i × HarvestTime_i + \\
& f2(Bathymetry_i) + b_{ID[i]} + f3(x_i, y_i)
\end{split}
\]

where:

\begin{itemize}
\tightlist
\item
  \(f1(·)\) and \(f2(·)\) are univariate smooth functions (thin-plate or
  cubic regression splines) of scaled continuous covariates;
\item
  \(DistHarvest_i\) is the (scaled) distance to nearest harvest feature
  for observation i (or distance in metres, scaled);
\item
  \(Sex_i\) is a categorical indicator (e.g., Male/Female); β1 captures
  the interaction of distance to harvest by sex;
\item
  \(HarvestTime_i\) is a categorical variable indicating harvest period
  class (e.g., before / during / after the June harvest); β2 captures
  the interaction of distance to harvest with harvest period;
\item
  \(b_{ID[i]}\) is an individual-level random intercept (a random
  effect) for the individual to which observation i belongs;
\item
  \(f_3(x_i, y_i)\) is a two-dimensional smooth over projected
  coordinates (x,y) to account for residual spatial structure.
\end{itemize}

Smoothing choices: univariate smooths (f1 and f2) used modest basis
dimension (k = 5) to prevent overfitting and to reflect expectation of
relatively simple environmental responses. The 2D spatial smoother f3
used a basis dimension k chosen by a pragmatic heuristic (see Appendix)
that balances flexibility and computational tractability.

\subsection{Estimation and computational
details}\label{estimation-and-computational-details}

\subsubsection{Estimation engine}\label{estimation-engine}

mgcv::bam was used for efficient estimation of GAMs on large data sets.
Smoothing parameters were estimated via (f)REML (method = ``fREML'') to
achieve stable penalised smoothness estimation in the presence of random
effects.

\subsubsection{Discrete approximation:}\label{discrete-approximation}

when datasets were large the discrete = TRUE option was used to speed
computation without materially altering estimates.

\subsubsection{Weights:}\label{weights}

presence points had default weight 1. Background (available) points
sampled from stratified areas received area-based weights (area of
stratum divided by number of sampled points in stratum) to reflect
unequal sampling intensity; weights were normalised to mean 1 so that
the effective sample sizes remain interpretable.

\subsubsection{Random effect treatment:}\label{random-effect-treatment}

individual-level heterogeneity was modelled as a random intercept via a
smooth term of the form s(ID, bs = ``re''). This captures consistent
differences among individuals that would otherwise inflate Type I error
due to non-independence of repeated measures.

\subsubsection{Spatial smoother:}\label{spatial-smoother}

s(x, y) models residual spatial dependence after accounting for
covariates and random effects, reducing spatially structured residual
variance that could confound covariate effects.

\subsubsection{Reproducibility:}\label{reproducibility}

random sampling of available points used a fixed seed; all final runs
record session information (software versions) to the archived results.

\section{Hypothesis tests, effect summaries and partial
effects}\label{hypothesis-tests-effect-summaries-and-partial-effects}

Primary inferential objectives include assessing whether distance to
harvest (and its interaction with sex and harvest period) influences
habitat selection probability. We summarise effects using:

Parametric coefficients and their standard errors (for interaction terms
and linear components). Estimated degrees of freedom (edf) and
approximate significance of smooth terms. Partial effect plots for
DistHarvest: to visualise the effect of distance to harvest on predicted
probability, we computed population-level predictions (see Section 7) on
the link (logit) scale while excluding individual random effects and the
2D spatial smooth, and transformed these to probabilities using the
logistic inverse. Confidence intervals were obtained by using the
standard errors of the link-scale predictions and back-transforming.
Relative selection (RS): we also present relative selection curves
computed as

\[
RS(d) = exp( η(d) − mean( η(d) ) )
\]

within each sex × harvest-period group, where \(η(d)\) is the predicted
link value at distance \(d\) and the mean is taken over the range of d
for the specific sex × period group. RS provides a multiplicative
measure of relative selection compared to the group mean and aids
ecological interpretation.

\section{Prediction, mapping and
interpretation}\label{prediction-mapping-and-interpretation}

Population-level maps: to visualise population-level spatial patterns in
use probability, we predicted p (probability of use) across the
prediction grid for each sex and for focal harvest periods. When
generating maps that reflect population-level baselines, we excluded
individual random effects (b\_ID) from predictions so that surfaces
represent average responses rather than idiosyncratic individual
differences.

Masking to habitat: predicted rasters were masked to the kelp habitat
(tareskog) so that maps are only shown in available habitat.

Uncertainty: per-grid-cell standard errors (on the link-scale) were
computed where possible and used to generate confidence bands for
partial effect curves. Spatial maps can be accompanied by uncertainty
maps (e.g., standard error rasters) to aid interpretation.

\section{Model validation and sensitivity
analyses}\label{model-validation-and-sensitivity-analyses}

Validation approaches:

\subsection{Model diagnostics from
mgcv:}\label{model-diagnostics-from-mgcv}

review of estimated degrees of freedom, significance of smooth terms,
concurvity measures, and residual plots.

\subsection{Spatial/block
cross-validation:}\label{spatialblock-cross-validation}

where implemented, we partitioned data into spatial or individual-based
blocks to assess out-of-sample predictive performance and to reduce
optimistic bias due to spatial autocorrelation. Reported diagnostics
include AUC (area under the ROC curve) and calibration plots for
held-out data.

Sensitivity analyses included varying accessible-area definitions
(individual-specific vs group-level), the number of background points
per used point, and harvest-stratification targets. We also varied the
basis dimension (nBasis) of the 2D spatial smooth to check the
robustness of spatial patterns to smoother flexibility. Interpretation
guidance:

Coefficients and smooths should be interpreted in the context of the
accessible-area assumptions; choice of accessible area can influence
estimates of selection. Results are therefore presented with
descriptions of the accessibility assumptions used in each analysis.

Significant interactions (e.g., DistHarvest × sex) indicate that the
effect of proximity to harvest on selection differs between sexes; plots
and RS curves are provided to visualise these interactions.




\end{document}
